% SPDX-FileCopyrightText: 2026 Bart van de Lint
% SPDX-License-Identifier: AGPL-3.0-or-later
\section{What it means for shipping}
\label{sec:industry}

For a shipowner the comparison turns on the fuel. On VLSFO at today's price wind pays, but slowly. The best kite pays back in 1.5 to 5.1 years, the rotors in 4.1 to 15 and the wing sails in 6.8 to 23 or never (\cref{tab:payback}), and the kite reaches those years only at its goals, while the rotors already have a service record (\cref{sec:res_goals}). On e-ammonia every rig pays back within a few years on every class, and the best rig closes about \NumPremiumClosed\,\% of the cost premium of the e-fuel over VLSFO (\cref{fig:premium}). Choose the rig with the fuel in mind, because the rig sized for VLSFO is smaller than the one sized for e-ammonia (\cref{fig:class_sizes}).

How much of shipping does each rig suit? \Cref{fig:fleet_share} adds up the class study in dollars, if each rig reaches its goals. For each class and route it takes the yearly cost of carrying the class's cargo and gives it to the rig that lowers the cost per tonne-kilometre most. That cost is the cost per tonne-kilometre of the ship without a rig, all ship costs and cargo interest included, times the 2018 transport work of the IMO size bin that class stands for \citep{faber2020} (the bins do not follow class names, so the 42,000 and 64,000 tonne deadweight Handysize and Ultramax stand for the bins just below their own size), the same work on both fuels and shared equally over the six routes, as no source says how each class divides its work between them. Where a second rig saves within \NumFleetLevelMargin\ point of the best, the largest error of a saving (\cref{sec:res_robustness}), the cost counts as level and goes to neither; where the best saves less than \NumFleetLevelMargin\ point, no rig is fitted. On VLSFO, of the USD \NumFleetBillOil\ billion a year the rotors take \NumFleetShareOilRotors\,\% and the kite \NumFleetShareOilKite\,\%, \NumFleetShareOilLevel\,\% is level, \NumFleetShareOilNone\,\% goes to no rig and the wing sails take \NumFleetShareOilWings\,\%. On e-ammonia, of USD \NumFleetBillAmmonia\ billion, the rotors take \NumFleetShareAmmoniaRotors\,\% and the kite \NumFleetShareAmmoniaKite\,\%, \NumFleetShareAmmoniaLevel\,\% is level, \NumFleetShareAmmoniaNone\,\% goes to no rig and the wing sails take \NumFleetShareAmmoniaWings\,\%. Together the rigs would save USD \NumFleetSavedOil\ billion a year on VLSFO and \NumFleetSavedAmmonia\ billion on e-ammonia. The thirteen classes carry \NumFleetWorkShare\,\% of the transport work of the bulk, tanker, container, general cargo and gas fleets and the largest car carriers, and emit \NumFleetFuelShare\,\% of shipping's CO$_2$; bins with no class, such as Suezmax tankers and gas carriers, are left out rather than assumed to follow a neighbour, and every ship is taken on the study's Atlantic routes. The shares assume perfect foresight of the weather. The kite keeps the least of its cost saving on forecasts (\cref{sec:res_robustness}), so its share is the one most likely to shrink against the rotors. \tbd{The VLSFO split is not settled: another draft of this study with a revised router, in which each rig's cost per tonne-kilometre moves by less than 1\,\%, gives the kite 55\,\% and the rotors 11\,\%, so read it only together with the level band.}

\begin{figure*}[tbp]
\centering
\plot[\textwidth]{fleet_share}
\caption{Share of the yearly cost of carrying the cargo of the thirteen ship classes (a) and of the yearly cost the rigs save (b) that goes to each rig, if each ship class and route fits the rig that lowers its cost per tonne-kilometre most and each rig reaches its goals, on VLSFO and e-ammonia. Each segment gives its share and, where it fits, its yearly US dollars; the total of each bar is under its fuel. Hatched: two rigs within the error of the cost saving, so neither can be named. Weighted by the 2018 transport work of the IMO size bins the classes stand for \citep{faber2020}. With perfect foresight of the weather. \tbd{The VLSFO split is not settled, see the text.}}
\label{fig:fleet_share}
\end{figure*}

How the cost saving is taken depends on how the ship trades. The router sets speed by cost, so part of each cost saving is time (\cref{tab:split}). A liner ship on a fixed schedule keeps the time of the ship without a rig and takes its gain as fuel instead (\cref{tab:speed_terms}). A tramp operator could gain more than our fixed round trips show (\cref{sec:disc_common}).

For a kite maker the goals are a development target and a business problem at once. Our kite price is that of roughly the thousandth unit built (\cref{sec:disc_prices}). In a business case of our own (\cref{sec:disc_details_prices}), a maker that spends about USD~100 million on development and sells 200 systems of \SI{800}{m^2} a year breaks even after about thirteen years at today's canopy price and later at the price goal. The spare canopies, not the first sale, pay for it, so a cheaper or longer-lived canopy works against the maker, and we keep the market factor of 2 on the canopy at the price goal. The goals say what a shipowner gains if they are reached, not that a maker earns back the development at this price. With every kite hardware and canopy price and the market factor at their upper bounds, close to the prices paid so far, the better of an \SI{800}{m^2} and a \SI{3200}{m^2} kite keeps \tbd{about half} of its cost saving on the \SI{140}{m} and \SI{200}{m} e-ammonia ships. On VLSFO it \tbd{costs more than it saves} (\cref{app:price_stage}). Wherever the kite beats the rotor in \cref{sec:res_routes}, it does so only in a mature market.

\tbd{A developer's work pays where \cref{fig:sensitivity} shows the largest gain: for a kite, a cheaper and longer-lived canopy, and for wing sails and rotors, a lower price per area and then more lift.} The numbers to measure first are those the comparison rests on and no one has measured: a ship kite's canopy life and its losses per launch, which decide whether the kite leads or trails the rotors (\cref{sec:res_missed}), and the share of time each rig is out of service at sea (\cref{sec:disc_reliability}).
